\section{Related work}
\label{sec:related}

This work sits at the meeting point of three lines: real-time optimisation under
uncertainty, conformal prediction and its conditional-validity theory, and the recent use
of conformal prediction inside optimisation. We review each in turn and then state the gap.

\subsection{Real-time optimisation under uncertainty}

The dominant paradigm for steady-state RTO under plant-model mismatch is modifier adaptation
\citep{marchetti2009,chachuat2009}, which adds first-order correction terms to the model cost
and constraints so that, on convergence, the model-based optimum is a KKT point of the plant.
Modifier adaptation corrects the nominal prediction, but it does not by itself supply a
probabilistic margin against an unmeasured disturbance. Safety against violation is handled
implicitly, through conservative tuning or through filtering of the modifiers.

A data-driven branch replaces or augments the modifiers with learned surrogates.
\citet{delrio2021}, extending the Gaussian-process modifier-adaptation line
\citep{delrio2019}, use Gaussian-process regression to capture the cost and constraint
mismatch non-parametrically, and they embed trust-region and acquisition-function ideas from
Bayesian optimisation to manage risk during exploration. There the GP mean corrects the
constraint and the GP variance drives excitation and sizes the trust region. The scheme is
probabilistic in its exploration, yet it does not issue a distribution-free, finite-sample
coverage guarantee on the constraint, and its risk control is trust-region containment rather
than a calibrated back-off. The closely related stochastic-MPC line
\citep{bradford2019,bradford2020} does subtract explicit GP-posterior back-offs from
constraints, but in a receding-horizon control setting and under a Gaussian posterior
assumption.

The classical route to explicit probabilistic constraints is chance-constrained programming
\citep{li2008}, which constrains the probability of satisfaction directly, computes
joint-constraint probabilities by sampling under an assumed, typically multivariate-normal,
disturbance law, and trades feasibility against profitability through the miss level. This is
the formulation we adopt in Eq.~\eqref{eq:chance}. The difference here is that the back-off
enforcing it is built distribution-free from data rather than from an assumed Gaussian law.

\subsection{Conformal prediction and conditional validity}

Conformal prediction \citep{vovk2005,angelopoulos2023} converts any point predictor into a
set with a finite-sample, distribution-free marginal coverage guarantee under
exchangeability. Split conformal and its locally weighted variants \citep{lei2018} rescale
the residual by a spread estimate so that the interval width adapts to the input, and
conformalised quantile regression \citep{romano2019} goes further by conformalising a
quantile-regression estimate so that the interval targets the conditional quantile while
retaining the marginal guarantee. All of these methods deliver coverage on average over the
calibration distribution.

The limit of that guarantee is the subject of a conditional-validity literature. Exact
distribution-free conditional coverage is impossible without further assumptions
\citep{vovk2012,leiwasserman2014}, and recent work characterises and partially restores
conditional guarantees \citep{gibbscherian2023}. A parallel strand relaxes exchangeability
itself, for distribution shift and dependent data
\citep{tibshirani2019,barber2023,gibbs2021}. The gap between marginal and conditional
validity is precisely what an optimiser exploits in this paper, because a margin calibrated
on average under-covers at the specific operating point an optimiser selects.

\subsection{Conformal prediction inside optimisation}

Conformal prediction has recently been brought inside chance-constrained optimisation.
Conformal predictive programming (CPP; \citealt{zhao2024cpp}) uses the conformal quantile
lemma to reformulate a chance-constrained program as a deterministic quantile-constrained
one, supporting any conformal variant and providing feasibility guarantees. Of central
importance for the present work, CPP observes that once a decision is optimised against
sampled constraint evaluations those evaluations are no longer independent of the decision,
so the exchangeability that underwrites the conformal guarantee is broken by the optimisation
itself, and it restores a guarantee a posteriori by recalibrating the selected solution on a
second, independent dataset. Related uses of conformal sets to certify safety in optimisation
and control include online-conformal safety guarantees for Bayesian optimisation
\citep{zhang2024} and conformal back-offs proposed for predictive control.

CPP supplies the a-posteriori calibration machinery this paper uses
(Section~\ref{sec:methods-apost}) and names the loss-of-independence mechanism in the
abstract. What it does not do, and what no work in the RTO literature above does, is
instantiate the mechanism in a process real-time optimisation loop, quantify how severely a
marginally valid, including locally adaptive, back-off is degraded by the optimiser, or
separate the two remedies the problem actually needs: a conditional construction to remove
the systematic selection bias, with a-posteriori calibration left to cover only the
finite-sample residual. CPP treats the a-posteriori step as the primary guarantee. We show
that without conditional validity that step is either large or infeasible, and that a
conditional back-off makes it small.

\subsection{Gap}

No existing treatment does three things together. It does not place a distribution-free
conformal constraint back-off inside a chemical-process RTO loop; it does not expose and
quantify the optimiser-induced selection effect that drives a marginally valid margin to
several times its nominal violation; and it does not identify conditional validity as the
structural remedy, relegating a-posteriori calibration to a residual correction and reading
its inflation factor as a diagnostic of data adequacy. This paper provides that treatment on
a rigorous debutaniser twin, and it maps the disturbance regime over which the treatment
holds. The soft-sensor framework of the companion paper \citep{companion2026} supplies the
model-based estimate of the unmeasured constraint that the back-off corrects; here that
estimate is taken as given and the focus is the back-off itself.
